\documentclass[11pt,a4paper]{article}
\usepackage[T1]{fontenc}
\usepackage{lmodern,amsmath,amssymb,amsthm,mathtools,booktabs,microtype,xcolor}
\usepackage[margin=26mm]{geometry}
\usepackage[colorlinks=true,linkcolor=blue!45!black,citecolor=blue!45!black,urlcolor=blue!45!black]{hyperref}
\newtheorem{theorem}{Theorem}[section]
\newtheorem{lemma}[theorem]{Lemma}
\newtheorem{proposition}[theorem]{Proposition}
\theoremstyle{remark}\newtheorem{remark}[theorem]{Remark}
\newcommand{\Q}{\mathbb Q}\newcommand{\Z}{\mathbb Z}\newcommand{\HH}{\mathcal H}
\setlength{\emergencystretch}{2em}
\title{An explicit quartic in 351 variables}
\author{}\date{9 November 2026}
\begin{document}\maketitle
\begin{abstract}
We specify an integer polynomial of degree four in 351 total variables
whose rational input projection is the integers, under the arithmetic
premises of the elliptic source construction. Its complete coefficient
table has 4990 nonzero monomials. Three height detectors use sixteen
local tests expressed by polynomials whose summands have different
valuation classes. Existing formal-logarithm coordinates and nine
additional field squares supply the required powers. The construction retains
the effective abscissa-sum height bound and the explicit fixed parameter
$80989^{54929}$. The polynomial has content one and a largest coefficient
with 1791357 binary digits.
\end{abstract}

\section{Statement and coefficient specification}
The rational input projection of a polynomial $P\in\Z[X,\boldsymbol Y]$
is the set of $x\in\Q$ for which $P(x,\boldsymbol y)=0$ has a solution
in rational values of $\boldsymbol Y$. Variable counts include the input
$X$, and degree always means total degree. The coefficient height $H(P)$
is the largest absolute integer coefficient.

For an integer $N\ge1$ put $A_N=80989^{2N+1}$. The exact tables
\nolinkurl{quadratic-family351.json} and
\nolinkurl{polynomial-family351.json} specify 323 primitive integer
quadratic rows $Q_{N,j}$ in 351 coordinates and
\[
                    P_N=\sum_{j=1}^{323}Q_{N,j}^2.
\]
The powers of $A_N$ are recorded separately in the tables solely to
abbreviate integer coefficients. The parameter is not a variable.

\begin{theorem}\label{thm:main}
Under the arithmetic premises in Section~\ref{sec:premises}, for every
integer $N\ge27464$ and every $x\in\Q$,
\[
 x\in\Z\quad\Longleftrightarrow\quad
 \exists\boldsymbol y\in\Q^{350},\ P_N(x,\boldsymbol y)=0.
\]
Every such polynomial has degree exactly four, content one, and 4990
nonzero coordinate monomials. Its coefficient height satisfies
\[
                         H(P_N)<2^{1520}80989^{4N+2}.
\]
In particular, substituting $A=80989^{54929}$ specifies a single
polynomial $P_{27464}$ with the asserted projection. Its largest
coefficient has exactly 1791357 binary digits.
\end{theorem}

\section{Arithmetic premises}\label{sec:premises}
Put
\[
 K=\Q(\rho),\quad\rho^2=2,\quad d=3-\rho,\quad
 E:y^2=x^3-d^2x,
\]
\[
 R_0=(1+\rho,2\rho),\qquad
 P=2R_0=\left(\frac{57-28\rho}{8},\frac{-660+417\rho}{32}\right).
\]
We retain the following arithmetic contract, stated with its exact
source locations in \cite[Section~2]{Source}.
\begin{enumerate}
\item The group and index-height assertions are
\[
 2E(K)=\Z P,\qquad h(x([n]P))\ge n^2-6\quad(n\ne0).
\]
Here $h$ is absolute logarithmic Weil height, and $\sigma$ denotes
conjugation on $K/\Q$.
\item For $t=u/v$ in lowest terms, $v>0$, outside
$\{0,1,2401,-343/32\}$, let $M_5(t)$ be the product of primes
$p>76176$ with a positive odd valuation on at least one of
\[
             v,\quad u,\quad u-v,\quad u-2401v,\quad32u+343v.
\]
The contact-height input is
\[
 h(t)\le\HH M_5(t)^{19},\qquad
 \HH=2^{110}210^{1800}\left(\prod_{p\le76176}p\right)^{19}+128
       <2^{2093448}.
\]
\item The source pole equations exclude negative odd valuations of their
nonzero input $b$ at primes outside $\{2,3,17\}$. They admit the specified
specialization when $0<b<17$, its reduced denominator is a square,
$b$ is a square in $\Q_2,\Q_3,\Q_{17}$, and its valuations there are
at least $3,1,2$, respectively. Their exact normalized seven-coordinate
form is retained in the coefficient table.
\item The simultaneous prime-pattern and conjugate reduction-order input
holds with arbitrary fixed odd avoidance modulus, the prescribed real
approximation, and the finite congruences simultaneously. The selected
primes are $1$ modulo $8\cdot3\cdot17$, with
$v_2(d_-)=6$, $v_2(d_+)=2$ and
$\gcd(d_-(p),d_+(p'))=4$. Its quantitative weighted form is retained:
for fixed distinct coefficients $c_j$, admissible patterns
$p_j=u+c_jw^2$, $N'/2\le u\le N'$,
$1\le w\le\tau(N')\sqrt{N'}$, have total weight at least
$c_0(N')^{3/2}\tau(N')$, individual weights
$O((\log N')^k)$, and $\tau(N')$ may decrease arbitrarily slowly
after the fixed avoidance and congruence conditions are imposed.
\end{enumerate}
These are arithmetic premises, not conclusions of coefficient
verification. The effective sum-height theorem used below is proved
in \cite{Sums}; it is not an extra conjectural premise.

\section{The local tests and the quadratic system}
The four root charts are
\[
 x=h^2,\qquad j^2=x^2-d^2,
\]
and give nonzero integer multiples $Q_i$ of $P$. The squareclass
argument establishing this assertion, including its converse and all
torsion exclusions, is proved in \cite[Section~2]{Baseline}.
Write their labels as $n,a,t,b$ and set
\[
 B=\max(|n|,|a|,|t|,|b|),\qquad
 S=\sum_{i=1}^4x(Q_i)=S_0+S_1\rho.
\]
The three positive rational square scales are
\[
                         1,\qquad S_0^2,\qquad(S_0+S_1)^2.
\]
They are nonzero because $S_0>12$ and $S_0>\sqrt2|S_1|$.

For each detector the contact input is $t_*=a_*\lambda$ for one of
these scales. Put $\delta=t_*-1$,
\[
 F_*(t_*)=t_*(t_*+1/2)(t_*+1)(t_*-1/2399)(t_*+32/407),
\]
\[
 R_*=\frac{F_*(t_*)}{a_*^5\delta^6},\qquad
 q_m=R_*^{1-2m}\delta^{-2m-2}.
\]
These rational formulas are used only when their denominators are
nonzero; the actual system consists of the polynomial rows specified
in the table. The exceptional value $t_*=1$ has empty contact radical
and is explicitly covered by the local proof in \cite{Packing}.

Write $f(Q)=1/x(Q)$, extended by zero at the identity. The auxiliary
points satisfy $U_0=H$, $V_0=H-P$, and $Q_i=U_i-V_i$ as group
differences on $E$. The exact addition rows and identity-inclusive
charts of \cite{Baseline} enforce these relations.

Define $A_i=f(U_i)$, $B_i=\sigma f(V_i)$, and
\[
 \begin{aligned}
 C_i&=L_i-L'_i\quad(1\le i\le4),& I&=L_1-\beta,\\
 D_H&=f(H),&G&=L_3^2-L_2(L_2-2)L_3L_4+(L_2-2)^2L_4^2-1.
 \end{aligned}
\]
The alternate ratios $L'_i$ already include conjugation. The nine tests
$A_i,B_i$ have order one, $C_1,I$ have order ten, and $D_H,C_2,C_3,C_4,G$
have order forty. For an ordered list of tests, use the factor
\[
 \mathcal T(D_0,\ldots,D_{s-1})
 =1+\sum_{j=0}^{s-1}q_{m(D_j)}^{2^j}D_j^{2^{j+1}},
\]
where $m(D_j)$ is its stated order. Thus different required orders can
occur within one factor.

The detector uses the product of the factors for these five lists:
\[
 \begin{gathered}
 (D_H,C_1,B_1,B_0),\quad(C_2,I,B_2),\quad(C_3,A_2,B_3),\\
 (C_4,A_3,B_4),\quad(G,A_4,A_1).
 \end{gathered}
\]
Each test occurs exactly once. The additional field powers are
\[
 C_1^4,I^4,A_2^4,A_3^4,A_4^4,\quad B_1^8,B_0^{16},\quad A_1^4,A_1^8.
\]
In particular the sixteenth power of $B_0$ is obtained by squaring its
existing eighth power, and the eighth power of $B_1$ by squaring its
existing fourth power. All additional field powers use quadratic rows.
The full normalization formulas and exact coordinate accounting are
in \cite{Packing}.

At a contact prime, a nonzero term in position $j$ has valuation
$2^j$ times an odd integer. These valuations are pairwise distinct
and different from zero. The unique-minimum argument and the literal
trace-clearing rows give every assigned test order at the same place;
this is \cite[Proposition~1.2]{Packing}. The argument also proves the
nonvanishing needed for the trace quotient. The rational logarithms,
point charts, and scalar equations retain their stated forms.

\section{Soundness with the explicit parameter}
The input label and the scalar normalizer are
\[
 \begin{gathered}
 \beta=51X^2+2504,\qquad
 r=\beta(\beta-1)(\beta-2401)(32\beta+343)(2\beta-1),\\
 R_s=\frac{A_Nr}{1+80989^2r},\qquad c_s=R_s+9792.
 \end{gathered}
\]
The retained scalar argument gives, with $c=|a-1|$,
\[
 c\ge80989^{L_N},\qquad
 L_N=\left\lfloor\frac{187(2N-1)}8\right\rfloor+1,\qquad
 h(\beta)\le\HH c^{19/24}.
\]
Its exceptional-point exclusions precede every use of the height
detectors and are unaffected by their polynomial factors. The integer
certificates in \cite{Baseline} give, for $N\ge27464$,
\begin{equation}\label{eq:margins}
 |a|>(40\HH)^{10},\qquad
 \log(48\HH)<|a|/40,\qquad h(\beta)<|a|/40.
\end{equation}

The effective sum-height theorem yields
\[
 B^2\le2h(S_0)+2h(S_0+S_1)+82.
\]
Combining it with the contact-height input for the three scales gives:
if a nonzero integer $J$ is divisible by the $k$th powers of all three
contact radicals, then
\begin{equation}\label{eq:height}
                          B^2\le5\HH |J|^{19/k}.
\end{equation}
The formal logarithm comparison uses only the orders $1,10,40$,
which the displayed factors supply unchanged. Therefore the fortieth
powers of the contact radicals divide
\[
 \Delta=t^2-a(a-2)tb+(a-2)^2b^2-1.
\]
If $n-\beta\ne0$, their tenth powers divide its nonzero reduced
numerator $J$.

\begin{proof}[Proof of soundness in Theorem~\ref{thm:main}]
A rational zero of $P_N$ makes all its quadratic rows zero.
If $\Delta\ne0$, then $|\Delta|\le14B^4$, and \eqref{eq:height}
with $k=40$ gives
\[
 B^2\le5\HH(14B^4)^{19/40}<20\HH B^{19/10}.
\]
Thus $B^{1/10}<20\HH$, contradicting $B\ge|a|$ and
\eqref{eq:margins}. Hence $\Delta=0$.

The integral Lucas descent in \cite[Section~5]{Source} applies to all
nonzero integer $b$ and $|a|>7$, giving $\log B>|a|-3$.
If $n-\beta\ne0$, its reduced numerator satisfies
$|J|\le2B\exp(h(\beta))$. Equation~\eqref{eq:height} with $k=10$
would imply
\[
 \log B<10\log(20\HH)+19h(\beta)<29|a|/40,
\]
contrary to the Lucas bound. Thus $n=\beta\in\Z$. Write $X=u/v$
in lowest terms with $v>0$. Integrality of $51X^2+2504$ implies
$v^2\mid51$. Since $51$ is squarefree, $v=1$ and $X\in\Z$.
\end{proof}

\section{Completeness and finite verification}
For an integer input, the scalar preparation, Lucas labels, and root
charts are those of \cite{Baseline}. They admit all nonzero integer
indices. The point and prime choices precede the fresh pole witnesses.
Proposition~3.1 of \cite{Packing} proves their simultaneous extension
to the present factors: at each selected contact prime the trace is
$1$ or $2$ modulo $p>2$; at the fixed finite places the same arbitrary
preparation depths give the required squareclasses; and on the fixed
compact real choice regions all weighted factors tend uniformly to one
as $\lambda\to+\infty$. Thus the pole inputs lie in $(0,17)$ and meet
the unchanged pole-realization hypotheses. Equal height scales simply
reuse identical witnesses.

All rational coordinate formulas are defined on these completeness
witnesses and fill the 323 rows. This proves the reverse implication
in Theorem~\ref{thm:main}. The numerical condition $N\ge27464$ was used
only in the unchanged soundness inequalities \eqref{eq:margins};
completeness introduces no further numerical bound on $N$.

The exact reproduction commands are
\begin{verbatim}
python build_quartic_mixed.py
python check_quartic_mixed.py
\end{verbatim}
The independent program reconstructs the field normalizations and every
modified and retained row, and compares the entire expanded coefficient
table. It verifies 323 quadratic rows, 351 coordinates, 5005 symbolic
terms before the fixed parameter is collected, and 4990 ordinary
monomials afterwards. It checks degree four, content one, the uniform
coefficient-height bound, and the exact largest-coefficient bit length
at $A=80989^{54929}$. The arithmetic premises in
Section~\ref{sec:premises} remain mathematical inputs.

\begin{thebibliography}{9}
\bibitem{Packing}
\emph{Simultaneous valuation conditions and explicit quartics},
9 November 2026, accompanying proof and independent exact verification.
\href{../local_tests/mixed_valuation_tests_20261109_v11.pdf}{PDF}; source: \nolinkurl{mixed_valuation_tests_20261109_v11.tex}.
\bibitem{Baseline}
\emph{A quartic in 370 variables with specified integer coefficients},
8 November 2026, complete quadratic specification and effective height proof.
\href{quartic370_20261108_v10.pdf}{PDF}; source: \nolinkurl{quartic370_20261108_v10.tex}.
\bibitem{Source}
\emph{A forced prime condition and an explicit quartic in 392 variables},
8 November 2026, arithmetic input contract and scalar/Lucas arguments.
\href{quartic392_20261108_v8.pdf}{PDF}; source: \nolinkurl{quartic392_20261108_v8.tex}.
\bibitem{Sums}
\emph{Effective heights of sums of elliptic abscissas}, 8 November 2026,
effective four-term sum-height theorem and its application to two rational
linear forms.
\href{../effective_sums/effective_elliptic_sums_20261108_v10.pdf}{PDF}; source: \nolinkurl{effective_elliptic_sums_20261108_v10.tex}.
\bibitem{Glazer}
E.~Glazer, \emph{An explicit quartic defining the integers in the rationals},
Version~13, manuscript and coefficient data, commit
\nolinkurl{5bbd74f6a479776a332cbe036809d5dddae63720},
\url{https://github.com/elliotglazer/diophantine-z-q}.
\end{thebibliography}
\end{document}
